% Chapter 1

\chapter{Estado del arte} % Main chapter title
\label{sec:Estado Arte} % For referencing the chapter elsewhere, use

\subsection{El Desbalance de Clases como Barrera Metodológica Persistente}

El desequilibrio de clases representa un desafío fundamental para la validez de los modelos de análisis de sentimientos. Este problema se manifiesta de manera consistente en corpora ampliamente utilizados donde la clase mayoritaria habitualmente excede el 70\% de las muestras. En conjuntos de referencia como Stanford Sentiment Treebank 2, TASS-SEPLN y Coursera Reviews, la ausencia de mecanismos correctivos provoca caídas pronunciadas en el macro F1-score y degradación generalizada del rendimiento \cite{gopali2024applicability}.

Este fenómeno trasciende barreras lingüísticas y dominios específicos. La subrepresentación de clases minoritarias conduce sistemáticamente a predicciones inexactas y bajo rendimiento del modelo para opiniones poco frecuentes \cite{gopali2024applicability,whitehouse2023llm}. La problemática no se limita meramente a métricas de exactitud global. Por el contrario, compromete la utilidad práctica de los sistemas al enmascarar su incapacidad fundamental para identificar patrones en categorías minoritarias, afectando tanto la precisión como la aplicabilidad en escenarios de producción.

Durante más de una década, este desequilibrio ha constituido una barrera metodológica reconocida por la comunidad científica \cite{gopali2024applicability}. La optimización ingenua de la exactitud global induce a los modelos a converger hacia estrategias que favorecen la predicción de la clase dominante, relegando sistemáticamente las opiniones minoritarias. La literatura reciente ha cristalizado las soluciones propuestas en cuatro corrientes metodológicas principales, cada una con fortalezas y limitaciones específicas que fundamentan la necesidad de enfoques híbridos más sofisticados.

\subsection{Taxonomía de Enfoques para Desbalance de Clases en Análisis de Sentimientos}

La investigación contemporánea ha desarrollado cuatro paradigmas fundamentales para abordar el desbalance de clases. Cada paradigma opera en diferentes niveles de abstracción y presenta compensaciones (trade-offs) distintas entre costo computacional, naturalidad lingüística y efectividad empírica.

\subsection{Remuestreo Geométrico Clásico en Espacios de Características}

El remuestreo geométrico constituye el enfoque más establecido. Algoritmos como SMOTE (Synthetic Minority Over-sampling Technique) \cite{chawla2002smote} y ADASYN (Adaptive Synthetic Sampling) \cite{he2008adasyn} generan muestras sintéticas mediante interpolación lineal en espacios vectoriales. Estudios recientes han documentado mejoras significativas en F1-score cuando SMOTE se aplica en conjunto con clasificadores basados en BERT para tareas de clasificación multiclase desbalanceadas \cite{cloutier2023fine}. La aplicación de estas técnicas busca mejorar la exhaustividad (recall) de clases minoritarias sin comprometer excesivamente la precisión global \cite{chawla2002smote,he2008adasyn}.

Sin embargo, investigaciones específicas en clasificación de texto han revelado limitaciones fundamentales de este paradigma. Cuando SMOTE opera exclusivamente en espacios de embeddings textuales, los vectores interpolados resultantes pueden carecer de correspondencia directa con texto lingüísticamente coherente. Esta restricción limita su utilidad en aplicaciones que requieren inspección humana o retokenización para propósitos secundarios. Un estudio comparativo exhaustivo utilizando datos de hoteles demostró que, aunque SMOTE integrado con métodos de ensemble alcanzó 88\% de precisión \cite{putra2025optimizing}, la naturalidad lingüística de los ejemplos sintéticos permanece como preocupación no resuelta.

Análisis adicionales en clasificación textual revelaron que SMOTE aplicado sobre vectores MiniLM produjo mejoras del 17\% al 20\% en recall, particularmente beneficiosas para clases minoritarias \cite{taskiran2025comprehensive}. No obstante, estos resultados vienen acompañados de una salvedad crítica: los puntos interpolados en espacios de alta dimensionalidad no garantizan validez semántica. Las estrategias complementarias de submuestreo (undersampling), como el submuestreo aleatorio (RUS), aunque reducen el sesgo hacia la clase mayoritaria, lo hacen a expensas de descartar información potencialmente valiosa. Este compromiso ha motivado la exploración de enfoques de selección más estratégicos y técnicas de ensamble que intentan preservar información mientras mitigan el desbalance.

\subsection{Sobremuestreo en Espacios de Características Profundas}

El segundo paradigma opera directamente sobre representaciones contextualizadas generadas por modelos transformer. Este enfoque interpola vectores producidos por arquitecturas como BERT \cite{devlin2018bert} con el objetivo explícito de preservar información semántica durante el proceso de síntesis. El método reconoce que los embeddings contextualizados capturan relaciones semánticas más ricas que representaciones basadas en frecuencia de palabras o embeddings estáticos.

Investigaciones recientes han validado este principio mediante implementación de Synthetic Feature Augmentation, que aplica variantes de SMOTE (incluyendo Borderline-SMOTE y ADASYN) directamente en espacios de embeddings BERT sin decodificación a texto \cite{synthetic2025feature}. Los resultados en clasificación binaria (IMDB, SST-2) y multiclase (AG News con cuatro categorías) demostraron que modelos aumentados con SMOTE superaron consistentemente baselines desbalanceados. Las mejoras fueron particularmente pronunciadas cuando el tamaño de la clase minoritaria es extremadamente reducido.

Visualizaciones mediante UMAP y t-SNE confirmaron que muestras sintéticas efectivamente llenan regiones de clase minoritaria manteniendo estructura de cluster y coherencia semántica. La contribución conceptual clave de este enfoque radica en demostrar que técnicas tradicionales de SMOTE, cuando se aplican a embeddings semánticos aprendidos en lugar de características crudas, permanecen altamente efectivas para tareas de procesamiento de lenguaje natural. A diferencia de metodologías que decodifican vectores interpolados a texto, este paradigma opera exclusivamente en espacio de embeddings, ofreciendo eficiencia computacional superior pero sacrificando interpretabilidad directa de las muestras sintéticas generadas.

\subsection{Generación de Texto Completo con Modelos de Lenguaje de Gran Escala}

El tercer paradigma aprovecha la capacidad de los Large Language Models para generar directamente texto completo, gramaticalmente correcto y contextualmente coherente que enriquece las clases minoritarias. Este enfoque ha demostrado ventajas sustanciales sobre métodos geométricos tradicionales en múltiples dimensiones.

La investigación seminal de Suhaeni y Yong documentó que la adición de reseñas sintéticas generadas por GPT-3 al corpus Coursera Reviews produjo incrementos del 12.7\% en exactitud media y hasta 13\% en macro F1-score \cite{suhaeni2023mitigating}. Este resultado estableció un benchmark crítico que motivó investigación subsecuente en múltiples contextos lingüísticos. En el dominio árabe, investigadores utilizaron AraGPT-2 para mejorar la detección de opiniones negativas en contenido de redes sociales \cite{habbat2023using}. Paralelamente, estudios multilingües demostraron mejoras de hasta 13.4\% en exactitud para tareas con datos de entrenamiento limitados, entrenando modelos más pequeños sobre datos sintetizados por GPT-4 \cite{whitehouse2023llm}.

Las comparaciones directas entre LLMs y remuestreo geométrico han revelado ventajas sistemáticas del primer enfoque. Gopali et al. alcanzaron F1-score de 0.76 utilizando GPT-3.5-Turbo, contrastando favorablemente con valores de 0.60 o inferiores obtenidos mediante combinaciones de SMOTE y Edited Nearest Neighbors en el corpus MBTI \cite{gopali2024applicability}. Investigación adicional de Cloutier y Japkowicz demostró que el sobremuestreo mediante LLMs generativos con fine-tuning (específicamente GPT-2) generalmente supera métodos de remuestreo tradicionales en clasificación multiclase desbalanceada \cite{cloutier2023fine}. Sin embargo, esta ventaja resulta menos pronunciada en clasificación binaria, y el fine-tuning del LLM emerge como paso crucial para la efectividad.

La calidad de datos generados constituye factor determinante del rendimiento. Evaluaciones humanas han destacado la capacidad de modelos como ChatGPT y GPT-4 para generar texto natural en múltiples lenguas, ocasionalmente superando la naturalidad de datos originales \cite{whitehouse2023llm}. No obstante, estas evaluaciones también señalan desafíos persistentes en consistencia lógica para ciertos idiomas de bajos recursos y mantenimiento de coherencia comparada con ejemplos producidos por humanos. Los LLMs de última generación capturan matices sutiles e identifican sentimientos mixtos con precisión superior a modelos clásicos, conservando eficiencia incluso en textos breves \cite{ghatora2024sentiment}.

A pesar de estos avances, este paradigma enfrenta limitaciones significativas. La investigación ha identificado que LLMs tienden a producir datos con mayor homogeneidad intra-clase, con incrementos del 14.51\% al 27.11\% en similaridad, potencialmente reduciendo la diversidad necesaria para entrenamiento efectivo \cite{howtogood2024synthetic}. Adicionalmente, la subjetividad tanto a nivel de tarea como de instancia se correlaciona negativamente con el rendimiento de modelos entrenados en datos sintéticos \cite{li2023synthetic}. Esto sugiere la necesidad de validación humana particularmente cuidadosa en dominios como análisis de sentimientos donde la ambigüedad interpretativa es inherente.

\subsection{Metodologías Híbridas Emergentes}

La convergencia de técnicas clásicas de oversampling con capacidades generativas de LLMs ha emergido como paradigma dominante en 2024 y 2025, respondiendo directamente a las limitaciones complementarias de enfoques previos. SMOTE identifica estratégicamente regiones subrepresentadas en espacios de embeddings pero carece de capacidad para generar texto lingüísticamente válido. En contraste, los LLMs producen texto coherente pero sin guía geométrica explícita sobre dónde posicionar muestras sintéticas. Esta complementariedad ha motivado el desarrollo de arquitecturas que utilizan SMOTE para identificación de posiciones estratégicas en el espacio de características, seguido de LLMs para materialización de estas posiciones como texto coherente.

Diversas implementaciones de este paradigma híbrido han explorado mecanismos de control y dominios específicos. La metodología LLMOverTab, aunque originalmente propuesta para datos tabulares, demostró que la combinación de generación textual mediante prompts con ajuste de densidad posterior usando SMOTE logra incrementos de 3\% en macro F1-score frente a la aplicación de cada técnica por separado \cite{isomura2025llmovertab}. Este resultado valida el principio de complementariedad entre generación y remuestreo geométrico.

SentiGEN, desarrollado por Sundararajan y Kumarapathirage, combina el modelo T5 con un algoritmo genético y un validador basado en XLNet \cite{sundararajan2023sentigen}. Esta arquitectura incorpora validación iterativa que rechaza muestras sintéticas de baja calidad, permitiendo que clasificadores entrenados solamente con datos sintéticos superen en rendimiento a aquellos entrenados con conjuntos de datos reales en corpus como SST-2 y Yelp. Esta capacidad aborda directamente la preocupación de control de calidad que limita otros enfoques.

AugmenToxic utiliza Aprendizaje por Refuerzo con Retroalimentación Humana para guiar la generación de LLMs hacia ejemplos de contenido tóxico difíciles de clasificar, elevando la exhaustividad en esta clase minoritaria específica \cite{bodaghi2024augmentoxic}. Paralelamente, LLM-AutoDA implementa un ciclo continuo de generación, fine-tuning y retroalimentación para descubrir la estrategia de aumento de datos óptima, superando a métodos estáticos en conjuntos de datos con distribuciones de cola larga \cite{wang2023llm}.

Las arquitecturas más representativas de este paradigma emergente establecen las bases conceptuales y metodológicas del campo \cite{smotext2025arxiv,imbllm2024arxiv}. La primera implementación documentada del paradigma híbrido SMOTE-LLM para texto ejecuta interpolación SMOTE estándar entre embeddings de alta dimensionalidad para identificar puntos sintéticos, seguido de decodificación mediante xRAG-7B que proyecta estos vectores interpolados a texto coherente \cite{smotext2025arxiv}. Los resultados en el dataset 20 Newsgroups demostraron que el enfoque alcanza rendimiento comparable a entrenamientos con datos reales. El hallazgo notable es que entrenar exclusivamente en datos sintéticos logró únicamente 3.2\% de degradación respecto al baseline con datos reales. Esta capacidad de aproximar distribuciones reales tiene implicaciones significativas para dominios con restricciones de privacidad.

Paralelamente, otras arquitecturas introdujeron estrategias de maximización de diversidad mediante condicionamiento en pares característica-valor aleatorios y fine-tuning del LLM exclusivamente en muestras minoritarias interpoladas, alcanzando F1-scores promedios superiores a 0.8 en datasets tabulares \cite{imbllm2024arxiv}. Un enfoque alternativo evita completamente la decodificación a texto, aplicando SMOTE y sus variantes directamente en espacios de embeddings BERT donde el clasificador opera \cite{synthetic2025feature}. Este paradigma sacrifica interpretabilidad humana de muestras sintéticas a cambio de eficiencia computacional superior, demostrando efectividad particular cuando el tamaño de clase minoritaria es extremadamente reducido. Las visualizaciones mediante UMAP confirmaron que las muestras sintéticas llenan efectivamente regiones minoritarias manteniendo coherencia semántica en el espacio latente, validando el principio de que técnicas tradicionales de SMOTE permanecen efectivas para procesamiento de lenguaje natural cuando se aplican a embeddings semánticos aprendidos en lugar de características crudas.

No obstante, estos enfoques híbridos presentan limitaciones metodológicas que motivan extensiones arquitecturales. La mayoría de arquitecturas aceptan outputs generados sin mecanismos de validación cuantitativa post-generación. Esto resulta problemático cuando LLMs producen texto fuera de tema, duplicados casi exactos de ejemplos de entrenamiento, o muestras que se desvían significativamente de la región pretendida en espacio de embeddings. Además, los métodos revisados aplican uniformemente peso unitario a todas las muestras sintéticas o utilizan ponderación global fija. Esta estrategia ignora que la confianza en la validez de una muestra sintética debería correlacionar con indicadores geométricos como distancia a ejemplos reales cercanos y densidad local de su región. Finalmente, algunos enfoques condicionan la generación únicamente en dos vecinos fijos identificados por interpolación SMOTE estándar, limitando la riqueza de información contextual disponible para el LLM, particularmente en clases con alta heterogeneidad intra-clase donde subclusters temáticos exhiben vocabulario y estructuras sintácticas distintas.

\subsection{Limitaciones Metodológicas y Oportunidades de Investigación}

El análisis de la literatura revela tres limitaciones fundamentales en los enfoques híbridos actuales. Estas limitaciones motivan el desarrollo de extensiones metodológicas más robustas.

La primera limitación se relaciona con el control de calidad post-generación. La mayoría de arquitecturas híbridas aceptan todos los outputs generados sin validación cuantitativa. Esta ausencia es problemática ya que los LLMs pueden producir texto fuera de tema, duplicados casi exactos de ejemplos existentes, o muestras que se desvían de la región pretendida en el espacio de embeddings. Los pocos métodos que implementan filtrado lo hacen mediante heurísticas sin análisis sistemático. No existe consenso sobre umbrales óptimos ni sobre cómo estos interactúan con características del dataset como ratio de desbalance o heterogeneidad intra-clase.

La segunda limitación concierne a la ponderación durante entrenamiento. Los métodos actuales aplican peso uniforme a todas las muestras sintéticas o utilizan ponderación global fija. Esta estrategia ignora información geométrica valiosa. Un sintético generado en región densa con múltiples vecinos próximos merece mayor confianza que uno posicionado en región dispersa. La confianza en la validez de una muestra sintética debería correlacionar con indicadores como distancia a ejemplos reales cercanos y densidad local de su región.

La tercera limitación se refiere a la selección de contexto para generación. Los métodos actuales condicionan la generación únicamente en dos vecinos fijos identificados por interpolación estándar. Esta restricción limita la información contextual disponible para el LLM. En clases con alta heterogeneidad intra-clase, donde existen subclusters temáticos con vocabulario y estructuras distintas, dos vecinos pueden ser insuficientes. Extender el conjunto de vecinos informativos mediante ponderación podría mejorar la coherencia y especificidad de clase de las muestras generadas.

Estas limitaciones son especialmente relevantes considerando los costos computacionales. Un pipeline que acepta todos los outputs sin filtrado genera costos significativos de API para producir muestras que pueden ser rechazadas durante entrenamiento por baja calidad. El desarrollo de mecanismos de filtrado adaptativo que mantengan tasas de aceptación entre 15\% y 20\% constituye una oportunidad clara de optimización. Solo las muestras de alta calidad deberían consumir recursos de entrenamiento.

Adicionalmente, la literatura carece de protocolos de evaluación estandarizados. Algunos estudios reportan únicamente exactitud global, otros macro F1-score, y pocos proporcionan análisis granular por clase minoritaria. La ausencia de estudios que aíslen contribuciones individuales impide la comprensión detallada de los mecanismos subyacentes. Esta carencia dificulta comparaciones objetivas y toma de decisiones informadas sobre qué técnicas adoptar para datasets específicos.

Finalmente, la extensión a tareas multiclase con desbalance extremo permanece insuficientemente explorada. La mayoría de estudios se enfocan en clasificación binaria o multiclase con 3 a 4 categorías. Datasets con 16 clases y ratios superiores a 20:1 representan desafíos de escalabilidad diferentes. El costo de generar sintéticos para todas las clases minoritarias puede volverse prohibitivo sin estrategias de priorización.

\subsection{Posicionamiento de la Presente Investigación}

La presente investigación aborda directamente las limitaciones identificadas mediante una arquitectura híbrida que integra cuatro innovaciones metodológicas. Primero, implementa una selección de soporte ponderada por K-nearest neighbors utilizando 8 vecinos contextualmente relevantes, ponderados por distancia inversa al punto sintético. Esta extensión proporciona información más rica para condicionar la generación del LLM, permitiendo capturar patrones locales más exhaustivos, particularmente en clases con heterogeneidad intra-clase significativa.

Segundo, la metodología introduce un sistema de filtrado multi-nivel adaptativo con tres filtros complementarios en cascada. El primer filtro valida similaridad con el vecindario local de la clase. El segundo previene duplicados casi exactos de datos de entrenamiento. El tercer filtro verifica que la muestra permanece en la región pretendida. Los umbrales pueden ajustarse dinámicamente según características de clase como tamaño y dificultad, permitiendo un balance entre cantidad y calidad de sintéticos generados.

Tercero, aplica ponderación basada en confianza geométrica durante el entrenamiento. Cada muestra sintética recibe peso proporcional a la inversa de su distancia promedio a vecinos de soporte. Las muestras posicionadas en regiones densas ejercen mayor influencia en el aprendizaje, mientras que aquellas en regiones dispersas, donde la confianza es menor, tienen peso reducido. Este mecanismo contrasta con enfoques que aplican peso uniforme independientemente del contexto geométrico.

Finalmente, el diseño enfatiza la eficiencia computacional sin fine-tuning. La arquitectura emplea generación directa con LLMs pre-entrenados mediante prompting contextualizado, sin requerir entrenar el modelo generativo exclusivamente en datos minoritarios ni implementar ciclos iterativos. Esta decisión reduce significativamente costos computacionales y tiempo de desarrollo. El filtrado adaptativo complementa esta eficiencia al asegurar que solo muestras de alta calidad sean aceptadas.

La validación empírica se estructura para abordar las brechas de evaluación identificadas. Mediante estudios de ablación sistemáticos, la investigación aísla contribuciones de cada componente. El análisis se extiende a tareas multiclase con desbalance extremo, específicamente datasets con 16 clases y ratios de hasta 27:1. Esto permite evaluar la escalabilidad del método en escenarios poco explorados en la literatura.

La contribución científica radica en la integración sistemática de estos componentes con justificación teórica y validación empírica rigurosa. El objetivo es proporcionar un protocolo reproducible que guíe la selección de técnicas según características del dataset y restricciones operacionales. Esto incluye ratio de desbalance, tamaño de clase minoritaria, heterogeneidad intra-clase, presupuesto computacional y requisitos de latencia.

